\chapter*{서문}
\addcontentsline{toc}{chapter}{서문}

이 책은 Hugging Face를 처음 만나는 실무자를 위한 실습형 원고입니다. 다만 시작점은 BERT가 아니라 훨씬 작고 투명한 도구인 \api{TfidfVectorizer}와 \api{LogisticRegression}입니다. 텍스트가 숫자가 되고, 숫자가 출력 헤드를 지나고, 출력이 Loss로 평가되는 흐름을 먼저 손에 익히기 위해서입니다.

현재 원고는 Part I, 1--6장까지를 담습니다. 각 장은 Google Colab 노트북을 원천으로 삼았고, 책에서는 실행 흐름을 최대한 유지하면서 조판과 설명의 리듬을 출판용으로 다듬었습니다.

\begin{noteBox}[읽는 방법]
코드를 먼저 따라 치기보다, 각 장의 변화추적표와 Loss 설명을 먼저 읽어 보세요. 그 다음 실습 코드를 실행하면 “이 장에서 정확히 무엇이 바뀌었는지”가 훨씬 선명하게 보입니다.
\end{noteBox}

\section*{박스 안내}

이 원고에는 몇 가지 반복되는 박스가 있습니다. 박스는 본문 흐름을 끊기 위한 장식이 아니라, 독자가 지금 어떤 종류의 정보를 보고 있는지 빠르게 구분하도록 돕는 장치입니다.

\begin{noteBox}[Chapter map]
각 장의 시작에는 원천 노트북, Colab 링크, 이 장의 초점이 정리됩니다. 책을 읽다가 직접 실행해 보고 싶을 때 이 박스를 기준으로 노트북을 열면 됩니다.
\end{noteBox}

\begin{previewBox}{미리보기}
미리보기 박스는 다음 장으로 넘어가기 전에 미리 알고 있으면 좋은 변화를 짚습니다. 토크나이저, Loss, Output Head가 다음 장에서 어떻게 달라지는지 가볍게 예고합니다.
\end{previewBox}

\begin{faqBox}{FAQ}
FAQ 박스는 학습 중 자주 생기는 실무적 질문과 이론적 질문을 함께 다룹니다. 질문과 답변을 하나의 시각 단위로 묶어, 필요한 부분만 다시 찾아 읽기 쉽도록 구성했습니다.
\end{faqBox}

코드 뒤에 붙는 \textbf{위 코드 읽기} 문단은 바로 앞 코드 블록의 핵심 줄을 풀어 설명합니다. 모든 줄을 기계적으로 해설하기보다, 모델 학습·토큰화·Loss 계산처럼 학습자가 놓치기 쉬운 의미 단위만 짚습니다.

\cleardoublepage
